\documentclass[a4paper,12pt]{ctexart}

\usepackage{tikz}% 画图用的包
\usepackage{graphicx} %插入图片的宏包
\usepackage{float} %设置图片浮动位置的宏包
\usepackage{subfigure} %插入多图时用子图显示的宏包
\usepackage{fancyhdr} %设置页眉页脚的宏包
\usepackage[]{caption2} %设置图和表的格式的宏包
\usepackage{multirow} %合并多行单元格的宏包
\usepackage{longtable} %不宽但很长的表格可以用longtable宏包来进行分页显示
\usepackage{array} %一般用于数学公式中对数组或矩阵的排版
\usepackage{makecell}% makecell命令对表格单元格中的数据进行一些变换的控制。我们可以使用 \ 命令进行换行，也可以使用p{(宽度)}选项控制列表的宽度
\usepackage{threeparttable} %制作三线表格
\usepackage{booktabs}%s三线表格中的上中下直线线型设置宏包，在这个包中水平线被教程\toprule、midrule、buttomrule。
\usepackage{enumerate} %列举宏包

% 以下是伪代码
\usepackage{algorithm}
\usepackage[noend]{algpseudocode}
\floatname{algorithm}{算法}


%页眉页脚设置
\pagenumbering{arabic}
\pagestyle{fancy}
\setlength{\headheight}{15pt}
\fancyhead[L]{\reportType}
\fancyhead[R]{\className \reportSemester}
\fancyhead[C]{}
\fancyfoot[C]{\thepage}

%标题序号长度设置
\setcounter{secnumdepth}{3}

%图片排版设置
\renewcommand{\figurename}{图} %重定义编号前缀词
\renewcommand{\captionlabeldelim}{.~} %重定义分隔符
 %\roman是罗马数字编号，\alph是默认的字母编号，\arabic是阿拉伯数字编号，可按需替换下一行的相应位置
\renewcommand{\thesubfigure}{(\roman{subfigure})}%此外，还可设置图编号显示格式，加括号或者不加括号
\makeatletter \renewcommand{\@thesubfigure}{\thesubfigure \space}%子图编号与名称的间隔设置
\renewcommand{\p@subfigure}{} \makeatother

%表头文字格式的详细设置
\renewcommand\theadset{\renewcommand\arraystretch{0.85}%
\setlength\extrarowheight{0pt}}%行距
\renewcommand\theadfont{\small}%字体
\renewcommand\theadalign{rt}%行列对齐
\renewcommand\theadgape{\Gape[0.5cm][2mm]}%上下垂直距离


\title{
  \begin{figure}[H]
    \centering
    \includegraphics[width=1\textwidth]{./img/university.png}
  \end{figure} 
  \vspace{3em}
  \huge \textbf{\reportName} \\ 
  \vspace{1em}
  \large \textbf{\className-\reportType}\\
  \vspace{5em}
  \large 学生姓名\hspace{0.7cm}\underline{\makebox[5.5cm]{\studentName}} \\
  \large 学生学号\hspace{0.7cm}\underline{\makebox[5.5cm]{\studentID}} \\
  \large 专业班级\hspace{0.7cm}\underline{\makebox[5.5cm]{\studentGrade}} \\
  \large 指导教师\hspace{0.7cm}\underline{\makebox[5.5cm]{\prof}} \\
  \large 提交日期\hspace{0.7cm}\underline{\makebox[5.5cm]{\the\year 年 \the\month 月 \the\day 日}} \\
}

\author{}
\date{}

\ctexset { section = { format={\Large \bfseries } } }



\newcommand{\studentName}{张三} % 学生姓名
\newcommand{\studentID}{20242100000} % 学生学号
\newcommand{\studentGrade}{计算机科学与技术 2024} % 学生班级

\newcommand{\reportName}{最小生成树算法比较}
\newcommand{\className}{算法分析与设计}
\newcommand{\reportType}{课程报告}
\newcommand{\reportSemester}{2024春}
\newcommand{\prof}{程振波}

\begin{document}
\maketitle
\newpage

\section{内容与要求}
本次课程报告以最小生成树问题为背景，请根据以下要求完成报告内容：
\begin{itemize}
    \item 问题背景请首先介绍什么是最小生成树问题，然后描述最小生成树问题的应用场景；
    \item 最小生成树算法介绍部分，请分别描述Prim算法和Kruskal算法，算法描述请使用伪代码。伪代码的格式可参考如下欧拉算法\ref{alg:euclid}的描述；
    \begin{algorithm}
        \caption{欧拉算法}\label{alg:euclid}
        \begin{algorithmic}[1]
        \Procedure{Euclid}{$a,b$}\Comment{a和b的公共}
        \State $r\gets a\bmod b$
        \While{$r\not=0$}\Comment{如果r是0}
        \State $a\gets b$
        \State $b\gets r$
        \State $r\gets a\bmod b$
        \EndWhile\label{euclidendwhile}
        \State \textbf{return} $b$\Comment{返回公共因子}
        \EndProcedure
        \end{algorithmic}
    \end{algorithm}
    \item 算法实现和比较部分包括：
    \begin{itemize}
        \item 要求分别使用不同数据结构实现Prim算法（数组和优先队列）和Kruskal算法（数组与并查集），并给出不同数据结构实现这两个算法运行时间的变化；
        \item 请给出选用不同数据结构对算法效率影响的理论分析；
        \item 比较Prim算法和Kruskal算法在求解稠密图和稀疏图时运行时间的比较；
        \item 算法比较需要随机生成至少10组不同的数据，然后统计每组不同数据的算法运行平均时间；
        \item 算法实现不要贴源代码，只需通过图表的方式给出不同情形（如数据结构）下的数据（如运行时间）；
    \end{itemize}
    \item 总结部分需要根据算法实现与比较的数据分析得出一般性的结论；
    \item 课程报告要求语言精炼，数据翔实，列出参考文献；
    \item 课程报告要求独立完成。
\end{itemize}

该课程报告的模版是在CTex\cite{ctex}基础上修改而成，如何在模版插入图片或者制作表格请参考该模版的文档。

\section{问题背景}

\section{最小生成树算法介绍}
\subsection{Prim算法}

\subsection{Kruskal算法}

\section{算法实现与比较}
\subsection{算法性能提升}

\subsection{算法适用场景分析}

\section{结论}
\begin{thebibliography}{9}
    \bibitem{ctex} https://ctex.org/
\end{thebibliography}
\end{document}